\begin{enumerate}[label=\thesubsection.\arabic*, ref=\thesubsection.\theenumi]
%---------------- Geometry of circles
\item Tangent--radius property. If $PA$ is a tangent to a circle with centre $O$ at the point $A$,
\begin{align}
	OA \perp PA \label{eq:CBSE25S2-tangent-radius}
\end{align}
\item Tangents from an external point. If $PA$ and $PB$ are tangents from $P$ to a circle with centre $O$,
\begin{align}
	PA = PB, \quad \triangle OAP \cong \triangle OBP, \quad \angle POA = \angle POB \label{eq:CBSE25S2-tangent-equal}
\end{align}
\solution
$OA = OB$ (radii), $PA = PB$ and $OP$ is common, so the triangles are congruent by SSS. 
\item Angle sum of a quadrilateral and of a triangle,
\begin{align}
	\sum \text{angles} = 360^\circ \text{ (quadrilateral)}, \quad 180^\circ \text{ (triangle)} \label{eq:CBSE25S2-angle-sum}
\end{align}
\item Angle in a semicircle. If $AB$ is a diameter and $C$ lies on the circle,
\begin{align}
	\angle ACB = 90^\circ \label{eq:CBSE25S2-semicircle}
\end{align}
\item Tangent--chord theorem. The angle between a tangent and a chord through the point of contact equals the angle in the alternate segment. If $PC$ is tangent at $C$,
\begin{align}
	\angle PCB = \angle CAB \label{eq:CBSE25S2-tangent-chord}
\end{align}
\item AA similarity. If two angles of one triangle equal two angles of another, the triangles are similar and corresponding sides are in the same ratio. For $\triangle LMN \sim \triangle ABC$,
\begin{align}
	\frac{LM}{AB} = \frac{MN}{BC} = \frac{LN}{AC} \label{eq:CBSE25S2-similar}
\end{align}
\item Section formula. The point $\vec{R}$ dividing $\vec{P}$ and $\vec{Q}$ internally in the ratio $m:n$,
\begin{align}
	\vec{R} = \frac{m\vec{Q} + n\vec{P}}{m+n} = \vec{P} + \frac{m}{m+n}\brak{\vec{Q}-\vec{P}} \label{eq:CBSE25S2-section}
\end{align}
For $m = n$ this is the mid-point $\frac{\vec{P}+\vec{Q}}{2}$.
\item Squared length and perpendicularity,
\begin{align}
	\norm{\vec{A}}^2 = \vec{A}^{\top}\vec{A}, \quad \vec{A} \perp \vec{B} \iff \vec{A}^{\top}\vec{B} = 0 \label{eq:CBSE25S2-dot}
\end{align}
\item Angle between two vectors. For direction vectors $\vec{m}_1$ and $\vec{m}_2$,
\begin{align}
	\cos\theta = \frac{\vec{m}_1^{\top}\vec{m}_2}{\norm{\vec{m}_1}\norm{\vec{m}_2}} \label{eq:CBSE25S2-angle}
\end{align}
\item Altitude to the hypotenuse. If $\angle BAC = 90^\circ$ and $AD \perp BC$, then
\begin{align}
	AD^2 = BD \cdot DC \label{eq:CBSE25S2-altitude}
\end{align}
\solution
Take $\vec{B} = \myvec{0\\0}$, $\vec{D} = \myvec{BD\\0}$, $\vec{C} = \myvec{BD+DC\\0}$, $\vec{A} = \myvec{BD\\h}$. From \eqref{eq:CBSE25S2-dot},
\begin{align}
	\brak{\vec{B}-\vec{A}}^{\top}\brak{\vec{C}-\vec{A}} = -BD\cdot DC + h^2 = 0
\end{align}
\item Zeroes of a quadratic. For $ax^2 + bx + c$ with zeroes $\alpha, \beta$,
\begin{align}
	\alpha + \beta = -\frac{b}{a}, \quad \alpha\beta = \frac{c}{a} \label{eq:CBSE25S2-zeroes}
\end{align}
\item Newton--Raphson method. To solve $f\brak{x} = 0$,
\begin{align}
	x_{n+1} = x_n - \frac{f\brak{x_n}}{f^{\prime}\brak{x_n}} \label{eq:CBSE25S2-newton}
\end{align}
\solution
The tangent to $f$ at $x_n$ is $y = f\brak{x_n} + f^{\prime}\brak{x_n}\brak{x - x_n}$. Setting $y = 0$ gives $x_{n+1}$.
\item Arithmetic progression with first term $a$ and common difference $d$,
\begin{align}
	a_n = a + \brak{n-1}d, \quad S_n = \frac{n}{2}\sbrak{2a + \brak{n-1}d} = \frac{n}{2}\brak{a + a_n} \label{eq:CBSE25S2-ap}
\end{align}
\item Gaussian elimination and rank. For a system $\vec{A}\vec{x} = \vec{b}$ with $n$ unknowns, row-reduce the augmented matrix $\myvec{\vec{A} & \vec{b}}$. The system is consistent if and only if
\begin{align}
	\rank\brak{\vec{A}} = \rank\myvec{\vec{A} & \vec{b}} \label{eq:CBSE25S2-rank}
\end{align}
and the solution is unique when this common rank equals $n$.
\item Simple interest on a principal $P$ at rate $r\%$ per annum for $t$ years,
\begin{align}
	I = \frac{Prt}{100} \label{eq:CBSE25S2-interest}
\end{align}
%---------------- Trigonometry
\item Standard values and inverse sine,
\begin{align}
	\sin 60^\circ = \frac{\sqrt{3}}{2}, \quad \tan 30^\circ = \frac{1}{\sqrt{3}}, \quad \tan 45^\circ = 1 \label{eq:CBSE25S2-trig}
\end{align}
The principal value $\sin^{-1}\brak{\frac{\sqrt{3}}{2}} = 60^\circ$.
\item Angle of depression. From $\vec{B}$ at height $h$ above the origin, the line of sight to a ground point with angle of depression $\theta$ has direction $\vec{m} = \myvec{\cos\theta \\ -\sin\theta}$ and meets the ground $\vec{e}_2^{\top}\vec{x} = 0$ at
\begin{align}
	\vec{B} + \lambda\vec{m}, \quad \lambda = -\frac{\vec{e}_2^{\top}\vec{B}}{\vec{e}_2^{\top}\vec{m}} = \frac{h}{\sin\theta} \label{eq:CBSE25S2-depression}
\end{align}
%---------------- Probability and statistics
\item Unit step, convolution and the $z$-transform. The unit step is $u\sbrak{n} = 1$ for $n \ge 0$ and $0$ otherwise, and
\begin{align}
	\brak{x * h}\sbrak{n} &= \sum_{k} x\sbrak{k}\,h\sbrak{n-k}, \quad
	X\brak{z} = \sum_{n} x\sbrak{n}z^{-n}
	\\
	\brak{x*h}\sbrak{n} &\xleftrightarrow{\ \mathcal{Z}\ } X\brak{z}H\brak{z} \label{eq:CBSE25S2-conv}
\end{align}
The running sum of a sequence is its convolution with $u\sbrak{n}$.
\item Probability by convolution. For a pmf $p\sbrak{n}$ and an event with indicator $m\sbrak{n} \in \cbrak{0,1}$, let $h\sbrak{n} = m\sbrak{-n}$. Then
\begin{align}
	\pr{E} = \sum_{n} p\sbrak{n}\,m\sbrak{n} = \brak{p * h}\sbrak{0} \label{eq:CBSE25S2-prob}
\end{align}
which is the coefficient of $z^{0}$ in $P\brak{z}H\brak{z}$. Also $\pr{E} + \pr{E^{\prime}} = 1$.
\item Grouped data. For a class with lower limit $l$, width $h$, frequency $f$, preceding cumulative frequency $cf$, and class mark $x_i$,
\begin{align}
	\bar{x} &= \frac{\sum f_i x_i}{\sum f_i} \label{eq:CBSE25S2-mean}
	\\
	\text{Median} &= l + \frac{\frac{N}{2} - cf}{f}\times h, \quad N = \sum f_i \label{eq:CBSE25S2-median}
\end{align}
The modal class is the class with the highest frequency. 
\end{enumerate}
